\section{The discrete energy and the common geometric limit}\label{sec:model}

This section records the exact relation between the walk and the electrical
network, and fixes the geometric representation used throughout the proof.
The representation retains every corner and every primal and dual incidence.
Consequently, a local electrical construction can later be compared with the
same limiting area measure and metric that occur in Theorem~\ref{thm:main}.

\subsection{The raw electrical energy}

For a finite map \(M\), choose an orientation of each unoriented edge only for
writing differences, and put
\[
 d_e f=f(e^+)-f(e^-),\qquad
 \mathcal E_M^0(f,g)=\sum_{e\in E(M)}d_e f\,d_e g.
\]
The value is independent of the chosen orientations.  Each parallel edge is
a separate summand, and a loop has zero difference.  We write
\(\mathcal E_n^0=\mathcal E_{M_n}^0\).  For a set \(F\) of edges,
\(\mathcal E_n^0(f;F)=\sum_{e\in F}(d_e f)^2\).

\begin{lemma}[Energy normalization]\label{lem:normalization}
Let \(Q_n\) be the unit-conductance graph Laplacian of \(M_n\), with loop
terms equal to zero, and set \(B_n=-2nL_n\).  On \(L^2(\mu_n)\),
\[
 B_n=\operatorname{diag}(\mu_n)^{-1}Q_n,\qquad
 \ip{f}{B_ng}_{L^2(\mu_n)}=\mathcal E_n^0(f,g).
\]
The kernel of \(B_n\) consists exactly of the constant functions.
In particular, its ordered eigenvalues are
\(\nu_{n,j}=2n\lambda_{n,j}\).
\end{lemma}

\begin{proof}
For a nonloop edge with endpoints \(u,v\), the two incident half-edge
contributions to \(-\ip{f}{L_ng}_{L^2(\mu_n)}\) are
\[
 \frac{1}{2n}\bigl(f(u)(g(u)-g(v))+f(v)(g(v)-g(u))\bigr)
 =\frac{d_e f\,d_e g}{2n}.
\]
The two half-edges of a loop contribute zero.  Summing proves both operator
identities.  A function of zero energy agrees across every nonloop edge, and
the map is connected, so it is constant.
\end{proof}

The superscript \(0\) will always indicate this unhalved electrical energy.
The continuum energy \(\mathcal E_h^0=2\mathcal E_h\) has operator \(B_h=2A_h\).
Thus conductivity one for the raw energies gives the clock in
Theorem~\ref{thm:main}: \(B_n/2=-nL_n\) converges to \(A_h\).

\subsection{Flags, corner occurrences, and the exploration}

We use the inventory encoding of the exact FK law in
\cite[Proposition~3.1 and Lemma~3.2]{C}.  At \(q=2\), its flexible-order
parameter is
\[
 b_{\mathrm F}=\frac{\sqrt2}{2+\sqrt2}=\sqrt2-1.
\]
The length-\(2n\) word conditioned to reduce to the empty word produces
precisely the probability law of Section~\ref{sec:introduction}.  Its resolved
contours have the parameter \(\gamma=\sqrt3\), and the limiting duration-one
surface is the ordinary unit-area quantum sphere
\cite[Theorem~3.4]{C}.  The convention \(m=2n\) therefore gives \(m\) letters
per unit of quantum-area time.

The finite exploration assigns a primal corner vertex \(v_n(i)\) to each cut
\(i\in\{0,\ldots,m-1\}\).  These are the \(m\) corner occurrences, including
repetitions of a vertex.  Consecutive cut vertices agree or are joined by a
primal edge.  We set \(v_n(m)=v_n(0)\) only to close the tour.  The repeated
closing cut receives no mass.  A vertex \(v\) occurs \(\deg(v)\) times, so
\begin{equation}\label{eq:corner-pushforward}
 \frac1m\sum_{i=0}^{m-1}\delta_{v_n(i)}=\mu_n.
\end{equation}
In particular, a vertex function is isometrically represented in
\(L^2([0,1],\dd t)\) by the function equal to \(f(v_n(i))\) on
\([i/m,(i+1)/m)\).  This representation will be used in the spectral
argument.

It is also useful to retain an actual surface on which both networks are
drawn.  For every occurrence of a vertex \(v\), an edge \(e\), and an incident
face \(F\), form an equilateral triangle with vertices labelled \((v,e,F)\).
Glue sides according to the cyclic incidences of the embedded map.  There are
\(4n\) triangles.  Repeated incidences are distinct triangles, even for a
loop or bridge.  The resulting incidence subdivision \(X_n\) is a topological
sphere with its piecewise Euclidean conformal structure.  Write \(m_n\) for
normalized triangle area; each flag has mass \(1/(4n)\).

If the interior of each flag is sent to its primal vertex, then
\begin{equation}\label{eq:flag-pushforward}
 (\pi_n)_*m_n=\mu_n.
\end{equation}
Indeed, \(v\) belongs to \(2\deg(v)\) flag occurrences.  This is also the
law of the tail of a uniformly chosen oriented root edge, conditional on the
unrooted decorated map.  The counting argument is valid in the presence of
automorphisms by counting rooted occurrences before taking the quotient
\cite[Lemma~2.2]{C}.

Choose three independent \(m_n\)-distributed points.  The unique
orientation-preserving uniformization
\(\phi_n:X_n\to\Sph\) sending them to \(0,1,\infty\) is the canonical
coordinate of \cite[Section~2]{C}.  We identify \(\Sph\) with the Riemann
sphere and write
\[
 z_n(v)=\phi_n(v),\qquad v\in V_n.
\]
The points used to normalize \(\phi_n\) are separate from exploration roots
and from any auxiliary chart marks.

\begin{lemma}[Duality in the flag surface]\label{lem:dual-geometry}
The map \((M,A)\mapsto(M^*,A^*)\), with a consistent directed dual-root
convention, is a bijection of the \(n\)-edge rooted decorated maps and
preserves their FK-Ising probabilities.  The primal and dual flag surfaces
are canonically isometric as unoriented piecewise Euclidean surfaces.  In
this identification a primal edge and its paired dual edge meet at their
common edge midpoint, with alternating incident rays; a primal edge meets no
other dual edge.  The same three physical flag-area samples may be used in
the two canonical uniformizations.  Their comparison is conformal or
anticonformal according to the orientation convention.
\end{lemma}

\begin{proof}
Taking complements across dual edges exchanges \(A\) and \(A^*\), and
\[
 \ell(M,A)=k(A)+k(A^*)-1=\ell(M^*,A^*).
\]
The dual operation is invertible, including its root convention, and keeps
the number of edges.  It therefore preserves the normalized finite law.
The flag correspondence is
\((v,e,F)\leftrightarrow(F,e^*,v)\).  It preserves the equilateral triangles
and their incidence gluings, including every repeated occurrence.  Drawing
a primal edge along its vertex-to-midpoint sides and a dual edge along its
face-to-midpoint sides proves the crossing assertion in each flag
neighborhood.  Conditional on the common metric surface, the same three
points are independent uniform area samples in both marginal experiments.
Genus-zero three-point normalization then gives the last assertion, allowing
complex conjugation if the surface orientation is reversed.
\end{proof}

\subsection{One coordinate for area, distance, and local tests}

We first record a normalization fact.  For a finite real random variable
\(Y\), recall that \(\mathfrak m(Y)\) is the zero of
\(r\mapsto\EE[\arctan(Y-r)]\).

\begin{lemma}[Continuity of the diameter normalization]\label{lem:space-normalization}
The center \(\mathfrak m\) is well defined, satisfies
\(\mathfrak m(Y+c)=\mathfrak m(Y)+c\), and is continuous under convergence
in distribution of finite real random variables.  If \(b_n>0\), \(b_n\to0\),
and \(b_n\Delta_n\) converges in distribution to the finite strictly positive
random variable \(\Delta\), then the normalization in
\eqref{eq:distance-normalization} satisfies \(a_n/b_n\to1\).
\end{lemma}

\begin{proof}
Bounded convergence shows that
\(F_Y(r)=\EE[\arctan(Y-r)]\) is continuous with limits
\(\pi/2\) and \(-\pi/2\) at the two ends of the real line.  It is strictly
decreasing because the integrand is strictly decreasing.  The zero exists
and is unique, and the shift identity follows directly.  If \(Y_j\) converges
in distribution to \(Y\), then \(F_{Y_j}(r)\to F_Y(r)\) for every fixed \(r\).
For each \(\varepsilon>0\), the two values at
\(\mathfrak m(Y)\pm\varepsilon\) have opposite strict signs; eventually the
same is true for \(F_{Y_j}\).  Its zero is between them, proving continuity.

Apply this fact to \(Y_n=\log(b_n\Delta_n)\).  The continuous mapping theorem
applies because \(\Delta>0\).  Shift equivariance gives
\[
 \log\frac{b_n}{a_n}
 =\mathfrak m\bigl(\log(b_n\Delta_n)\bigr)-\mathfrak m(\log\Delta)
 \longrightarrow0.
\]
\end{proof}

The following proposition collects the geometric results of the two
companions at the scope used here.  The canonical companion supplies the
uniformization and full flag geometry.  The metric-measure companion is used
in its revision of October 3, 2026.  Local conditional laws will be stated
separately in the next section.

For each rerooted contour, \(S_{\mathrm{ref}}\) denotes the sphere
reconstructed from its limit, written in a chosen conformal coordinate
with spherical metric \(\rho_{\mathrm{ref}}\), and
\(\eta^{\mathrm{ref}}:[0,1]\to S_{\mathrm{ref}}\) is its
area-parametrized tour.  Write \(P_n:X_n\to S_{\mathrm{ref}}\) for the
continuous contour projection and \(\Psi_n:X_n\to S_{\mathrm{ref}}\)
for the geometric homeomorphism below.  For a bilateral contour use the
one-point compactification of its reference plane and
\(\eta^{\mathrm{ref}}:\mathbb R\to
S_{\mathrm{ref}}\setminus\{\infty\}\); the local homeomorphism maps
the finite sewn flag sphere to that compactification.

\begin{proposition}[Common geometric extraction]\label{prop:geometry}
From every sequence \(n\to\infty\), one may pass to a further subsequence
and a representation coupling with the following almost-sure properties.
The limiting coordinate sphere carries an ordinary unit-area
\(\sqrt3\)-LQG sphere with area \(\mu_h\) and metric \(D_h\).
\begin{enumerate}[label=\textup{(\roman*)}]
\item The pushed corner measures converge weakly,
\[
 (z_n)_*\mu_n\longrightarrow\mu_h,
\]
and, in the product of the round topology with the ordinary topology of
\([0,\infty)\),
\begin{equation}\label{eq:common-distance-graph}
 \bigl\{(z_n(u),z_n(v),a_nd_n(u,v)):u,v\in V_n\bigr\}
 \longrightarrow
 \bigl\{(z,w,D_h(z,w)):z,w\in\Sph\bigr\}
\end{equation}
in Hausdorff distance.  Also
\(\delta_n:=\max_T\diam_{\mathrm{round}}\phi_n(T)\to0\), where \(T\) ranges
over flags.  Every complete primal or dual incidence star consequently has
diameter at most \(2\delta_n\).

\item Finitely many conditionally uniform rerootings may be retained using
the same map and the same \(\phi_n\).  For each rerooting \(r\), there is a
continuous tour \(\eta_r:[0,1]\to\Sph\) with
\((\eta_r)_*(\dd t)=\mu_h\) and
\begin{equation}\label{eq:uniform-tour-tracking}
 \sup_{0\le t\le1}
 d_{\mathrm{round}}\!\left(
 z_n\bigl(v_n^{(r)}(\lfloor 2nt\rfloor)\bigr),\eta_r(t)\right)
 \longrightarrow0.
\end{equation}
For \(\mu_h\)-almost every \(z\), the fiber
\(\eta_r^{-1}(\{z\})\) consists of one time among all of \([0,1]\).
For two conditionally independent uniform rerootings, their limiting root
positions, jointly with \(\mu_h\), have the law of two independent samples
from \(\mu_h\) conditional on \(\mu_h\).  They are distinct almost surely.

\item For each finite contour reference sphere, the projection \(P_n\)
places every point of each flag at a vanishing maximum
\(\rho_{\mathrm{ref}}\)-distance from \(\eta^{\mathrm{ref}}(t)\),
over all tour representatives \(t\) of its quadrangle corners and
all occurrences of its incident primal or dual vertices.  The
homeomorphisms satisfy
\[
 \sup_{x\in X_n}
 \rho_{\mathrm{ref}}\bigl(\Psi_n(x),P_n(x)\bigr)\longrightarrow0.
\]
On further subsequences,
\(\phi_n\circ\Psi_n^{-1}:S_{\mathrm{ref}}\to\Sph\) converges
uniformly in the target round metric on \(\Sph\) to the conformal or
anticonformal three-point comparison map.  The sign is retained in
the coupling.

For each member \(K\Subset U\) of a fixed finite family, protection
chooses a finite union \(I\) of tour intervals with
\[
 (\eta^{\mathrm{ref}})^{-1}(K)\subset I
 \Subset(\eta^{\mathrm{ref}})^{-1}(U).
\]
Eventually \(I\) contains every tour representative of each flag
whose contour projection meets \(K\) and of the full incidence stars
of its three vertices.  Intervals are interpreted on the time circle
for a sphere and are bounded in the bilateral case.  In a protected
bilateral window the finite sewing leaves these stars unchanged,
including cyclic order and original primal incidences.  The
homeomorphism of this protected bilateral sewing has projection error
tending to zero over the tests, and its full inverse image of \(K\)
lies in the unchanged part.  The eventual index may depend on the
tests and the sample.

\item For either rerooting with root \(p_*\), adjoin two independent
auxiliary pairs of \(\mu_h\)-samples conditional on the surface,
independently of retained exploration and chart observations given
that surface and its contour.  Use the separate three-marked
representations
\((Q_1,Q_2,p_*)=(0,1,\infty)\) and
\((Q_3,Q_4,p_*)=(0,1,\infty)\).
In either representation the chart interiors are relatively compact
disks \(U\Subset\mathbb C\setminus\{0,1\}\), with the required
strict buffers, on which the law of the actual-height sphere
real-field restriction to \(U\) is absolutely continuous with
respect to the local reference field law fixed for that comparison
in Section~\ref{sec:locality}.  The two chart families cover
\(\Sph\setminus\{p_*\}\); the families for two distinct roots have
a finite subcover of \(\Sph\).  Each comparison is used only in its
own three-marked marginal, without conditioning on the other auxiliary pair.
\end{enumerate}
The geometric conclusions persist on any further jointly tight extraction
that retains additional random observations of the maps and contours.
\end{proposition}

\begin{proof}
The canonical limit \cite[Theorem~1.1]{C} gives weak convergence of
\((\phi_n)_*m_n\), Hausdorff convergence of the full distance graph with a
deterministic scale \(b_n\to0\), and vanishing maximum flag diameter in the
same coordinate.  The third coordinate of the graph uses the uncompactified
distance, so its maximum converges to
\(\Delta=\diam(\Sph,D_h)\).  This diameter is finite and strictly positive.
Replacing the graph diameter by \(\Delta_n=\max\{1,\diam(V_n,d_n)\}\)
does not affect this limit because \(b_n\to0\).  Lemma~\ref{lem:space-normalization}
gives \(a_n/b_n\to1\), hence the graph convergence in
\eqref{eq:common-distance-graph} and \(a_n\to0\).  Couple each flag-area
point to its primal vertex.  Their images are within \(\delta_n\); therefore
\eqref{eq:flag-pushforward} gives the weak corner-measure convergence.
All flags in one incidence star contain the star's center, which proves the
\(2\delta_n\) bound.

The duration-one area tour and its continuity come from
\cite[Theorem~3.4]{C}.  The homeomorphisms of
\cite[Theorem~4.8]{C} track every flag and its tour representatives uniformly.
The uniformization comparison \cite[Theorem~8.8]{C} has only the conformal
and anticonformal three-point limits.  Passing further and retaining the
sign gives the uniform tracking in the canonical coordinate; uniform
continuity of the limiting tour extends the cut-time statement to
\eqref{eq:uniform-tour-tracking}.  These arguments apply jointly to a
finite number of rooted marginals by tightness and a further extraction.
Their pushed area measures are the same limit of the physical flag measures,
so each tour has pushforward \(\mu_h\).

Lemma~3.7 of \cite{C} says that an independent quantum-area point has exactly
one visit by the tour.  Integrating this probability-one statement
conditionally on the surface and tour gives the asserted singleton fibers
for \(\mu_h\)-almost every \(z\); a finite intersection handles all retained
rerootings.  Conditional uniform rerooting has tail law \(\mu_n\) by
\eqref{eq:flag-pushforward}.  Weak convergence of these measures, tested
against products of bounded continuous functions, identifies the joint law
by first taking \(\mu_h\) and then independently sampling the two roots
from it.  Quantum area is
nonatomic, so two such roots are distinct.

The continuous projections and complete time-preimage protection are
\cite[Proposition~4.2 and Lemma~4.3]{C}.  The reference homeomorphisms
and their validity with extra observations are exactly
\cite[Theorem~4.8]{C}; the canonical comparison is
\cite[Theorem~8.8]{C}.  The protected bilateral sewing and its inverse-image
capture are \cite[Proposition~4.9]{C}.  These statements concern fixed tests
with room and supply no rate for a test whose boundary changes with \(n\).

Finally, the sphere field comparison in \cite[Lemma~18.1]{P} uses the
exploration root and an auxiliary pair of area marks.  Use two separate
pair marginals.  Their four marks are distinct almost surely, and the
complements of the two pairs cover the sphere off the root.  Strict chart
interiors exhaust each complement.  Repeating this for two distinct roots
covers the entire sphere.  Compactness yields a finite subcover.  This
argument intersects probability-one conclusions from the separate
marginals; it does not strengthen either conditional field comparison.
\end{proof}

All subsequent subsequential arguments take place in such couplings, enlarged
by the compact local data they require.  A property proved almost surely on
every further extraction will be returned to convergence in distribution at
the end.  The choice of conformal or anticonformal sign has no effect on the
round Dirichlet integral.
